\section{Integer recourse, flows, and random components}\label{sec:int}

This section treats residual variables that are native integers. The core
search of \cref{sec:rec} applies unchanged, because it uses only the
conditional value function on a fixed residual set. Closure changes. A
patch and a Hessian test are meaningless for integer coordinates; instead,
a few bound-restricted recourse calls certify that one residual integer
label is optimal throughout the retained core hull, and an exact solver for
the small continuous core completes the problem. Under core-only noise
residual ties can persist on every draw, and for integer flows and totally
unimodular systems the label certificate is replaced by one that works with
the whole set of tied optimal labels. The section ends with two results of a
different kind: exact optimization of mixed boxes under strong noise through
random components, and an exact lattice argument for separable integer
problems with a low-rank concave part.

\Cref{tab:int:summary} lists the results. The factor $Q_{\rm ex}$ is defined
in \eqref{eq:rec:Q}, and $c_d$ is the constant of \cref{thm:int:solver}.

\begin{table}[t]
\centering
\caption{Integer recourse and random components. ``Label $+$ algebraic'' is
an integer residual label together with the algebraic output of
\cref{def:model:outputs}\ref{it:model:algebraic} for the core and the optimal
value; its size bound holds on every draw. The function $f_d(k)$ is
effective.}
\label{tab:int:summary}
\small
\begin{tabularx}{\textwidth}{@{}lXlll@{}}
\toprule
Result & Noise; residual structure & $\log_2M$ & Expected work & Output\\
\midrule
\cref{thm:int:native} & ambient; exact integer oracle on tightened bounds
& $\poly_d(I)$ & $[8^kQ_{\rm ex}+c_d^k]\poly_d(I)$ & label $+$ algebraic\\
\cref{cor:int:native-implicit} & ambient; as above & $\poly_d(I)$
& $8^kQ_{\rm ex}\poly_d(I)$ & label $+$ implicit\\
\cref{thm:int:flow-interior} & core-only; convex integer flow, interior
optimal cores & $\poly_d(I)$ & $[8^kQ_{\rm ex}+c_d^k]\poly_d(I)$
& label $+$ algebraic\\
\cref{thm:int:flow-boundary} & core-only; convex integer flow
& $f_d(k)\poly_d(I)$ & $f_d(k)Q_{\rm ex}\poly_d(I)$ & label $+$ algebraic\\
\cref{cor:int:bilinear} & core-only; flow, bilinear coupling & $\poly_d(I)$
& $[8^kQ_{\rm ex}+c_d^k]\poly_d(I)$ & label $+$ algebraic\\
\cref{thm:int:tu} & core-only; totally unimodular & as for flows
& as for flows & label $+$ algebraic\\
\cref{thm:int:strong-field} & ambient, strong; any mixed box & any
& see \eqref{eq:int:sf-work} & component\\
\cref{thm:int:lowrank} & aligned; convex piecewise-polynomial unaries on an
integer box & $\poly(I)$ & $8^kH_{\rm rat}\poly(I)$ & rational, no fallback\\
\bottomrule
\end{tabularx}
\end{table}

\subsection{An exact solver for a small continuous core}\label{sec:int:solver}

All label certificates end with the same task: minimize an explicit
polynomial exactly over a box of small dimension. A general exact method
returns separate defining polynomials for the coordinates; pairing their
roots costs a product of degrees, and it does not by itself say which
coordinate roots belong to the same optimizer. The following solver returns
all coordinates and the optimal value as polynomials in one common real
algebraic number, with a base that is exponential only in the dimension.

\begin{theorem}[Exact box solver through joint critical limits]\label{thm:int:solver}
Fix $d\ge1$, and let $D$ be the least even integer greater than $d$. There
are a constant $c_d$, depending only on $d$, and a deterministic algorithm
with the following properties. The input is an explicit rational polynomial
$f$ of degree at most $d$ in $k$ variables and a nonempty closed rational box
$B\subseteq\R^k$; let $H$ be its length. The algorithm returns a global
minimizer $x^*$ of $f$ on $B$ and the value $f(x^*)$ in the algebraic format
of \cref{def:model:outputs}\ref{it:model:algebraic}: one real root $\theta$
of a squarefree integer polynomial $P$ of degree at most $(D-1)^k$, given
with a rational isolating interval, and polynomials
$r_1,\ldots,r_k,r_f\in\Q[T]$ of degree less than $\deg P$ with
\[
 x^*_i=r_i(\theta)\quad(i=1,\ldots,k),\qquad
 f(x^*)=r_f(\theta),\qquad
 r_f=f(r_1,\ldots,r_k)\bmod P .
\]
Coordinates whose interval is a point have constant maps. As an auxiliary
representation, used only for comparisons, the algorithm also returns a
nonzero integer polynomial with the root $f(x^*)$ and an isolating interval
of that root. The bit work and the output length are at most
$c_d^k\poly_d(H)$. For every $q\ge0$, a rational point within Euclidean
distance $2^{-q}$ of $x^*$, a feasible rational point with certified
objective gap at most $2^{-q}$, and a rational interval of width at most
$2^{-q}$ containing $f(x^*)$ are computed within $c_d^k\poly_d(H+q)$. No
genericity is assumed: ties, singular critical points, and
positive-dimensional sets of minimizers are allowed.
\end{theorem}

For each face of the box, the solver deforms the restricted objective by
$\varepsilon\sum_ix_i^D$. Written in $z=1/\varepsilon$, the deformed
stationary equations have the pure-power leading monomials $x_i^{D-1}$ over
the field of rational functions in $z$ and after every finite specialization
of $z$. In particular this proves the claim for symbolic $\varepsilon$ and
for every $\varepsilon\ne0$.
They therefore form a Gr\"obner basis, and the quotient algebra has the
fixed dimension $(D-1)^r$ on a face with $r$ free coordinates. At
$\varepsilon=0$ the deformation disappears and the original stationary set
can be positive dimensional; the limit $\varepsilon\to0$ is handled by a
branch analysis, not by specializing. View the complex critical points of
the deformed problem as branches parametrized by $\varepsilon$. Call a linear
form $\lambda^{\mathsf T}x$ \emph{good} if it does not vanish on the leading
coefficient vector of any branch that escapes to infinity as
$\varepsilon\to0$, and if it separates the distinct finite limits of the
bounded branches. For a good form, the
leading coefficient in $z$ of the characteristic polynomial of
$\lambda^{\mathsf T}x$ in the quotient algebra vanishes exactly at the
projections of these finite limits, and its derivatives with respect to
independent coefficients of $\lambda$, divided by a greatest common divisor,
recover all coordinates of each limit as polynomials in its projection. For
other forms the leading coefficient can retain bounded projections of
escaping branches, so its roots need not be projections of limits. A
deterministic list of $O(r(D-1)^{2r})$ moment-curve forms contains a good
form. Every listed form is screened by degree and divisibility tests, and a
point reconstructed from any form is kept only after an exact test that it
lies in the box. A bad form that passes the screening can therefore add only
feasible candidates, whose values are at least the minimum. The deformed
minimizers have a convergent subsequence whose limit is a global minimizer of
$f$, and a good form recovers it, so the kept candidate of least value is
optimal. The proof is in \cref{app:int:solver}; it uses finite-dimensional
quotient algebras, simultaneous triangularization, Puiseux series, and exact
univariate algorithms, and \cref{lem:int:budget} gives the explicit choice
$c_d=2^{10000D}$ for one elementary implementation. That constant is a proved
bound used to define noise laws below, not a practical estimate. The
ingredients are classical. Recovering all coordinates from one univariate
polynomial and derivatives of a characteristic polynomial is the mechanism of
the rational univariate representation \citep{rouillier1999-solving-zero-dimensional-systems-through},
and deformations of stationary systems are standard in critical-point
methods \citep{basu2006-algorithms-in-real-algebraic-geometry}. We give a self-contained proof for boxes,
including singular and positive-dimensional critical sets, because the
recourse theorems need its exact contract and its constant-base bound; we do
not claim a new general algebraic complexity result.

\begin{corollary}[Mixed boxes and exact comparisons]\label{cor:int:mixed-solver}
For a mixed box with $k$ continuous coordinates and native integer
coordinates with $N_1,\ldots,N_t$ labels, enumerating the labels and applying
\cref{thm:int:solver} to each gives an exact global minimizer and value, in
the algebraic format with the integer labels listed, in
$c_d^k\prod_iN_i\poly_d(H)$ bit operations. Values of two algebraic outputs
of \cref{thm:int:solver} can be compared exactly, including equality, within
$c_d^k\poly_d(H)$; the comparison uses the two auxiliary value polynomials
and forms no field containing both.
\end{corollary}

\subsection{Native integer recourse}\label{sec:int:native}

\begin{definition}[Native integer recourse]\label{def:int:native}
A \emph{native integer recourse instance} is a core--recourse instance
(\cref{def:rec:instance}) whose residual set is
\[
 Y=\{z\in\Z^{n_R}:\ Az\le b,\ \ell\le z\le u\},
\]
with a rational matrix $A$, a rational vector $b$, and integer bounds
encoded in binary; the curvature bound \eqref{eq:rec:L} is required on
$[0,1]^k\times\prod_i[\ell_i,u_i]$ (verified by monomial bounds or
certified). Put $N_Z=\prod_i(u_i-\ell_i+1)$ and $S=k+\sum_i(u_i-\ell_i)$. A
\emph{box-stable exact integer oracle} receives a rational core point $v$,
rational linear coefficients, and integer bounds $\ell\le\ell'\le u'\le u$,
and returns the exact optimal value of
$\min\{F_\gamma(v,z):z\in Y,\ \ell'\le z\le u'\}$ and an attaining integer
point, or reports infeasibility, in bit work polynomial in the query length
with an exponent independent of $k$.
\end{definition}

The core changes costs, never feasibility: supplies, right-hand sides, and
bounds do not depend on $v$. The residual set is nonempty by
\cref{def:rec:instance}; one oracle call with the original bounds decides
this before sampling, and an empty $Y$ is reported as such. Exact values are
rational because the core point is rational and residual points are
integral. The oracle must handle arbitrary tightened coordinate bounds,
which is how competing labels are excluded.

\begin{theorem}[Native integer recourse]\label{thm:int:native}
Fix $d$. Let a native integer recourse instance with $k\ge1$ be given,
together with a box-stable exact integer oracle. There is a power of two $M$,
computed from the base instance with $\log_2M\le\poly_d(I)$, such that for the
ambient perturbation model with marginal law $U_{\sigma,M}$ the following
hold.
\begin{enumerate}[label=(\roman*)]
\item On every draw the algorithm returns a global minimizer $(v^*,z^*)$ of
$F_\gamma$ on $X=[0,1]^k\times Y$ and the optimal value: the label $z^*$
exactly, and $v^*$ and the value in the algebraic format of
\cref{thm:int:solver}. Its length is at most $c_d^k\poly_d(I)$ on every
draw, and a $q$-bit refinement costs $c_d^k\poly_d(I+q)$.
\item The expected bit work and the expected size of the global proof record
are at most $[8^kQ_{\rm ex}+c_d^k]\poly_d(I)$.
\item If $F_0$ is quadratic, the output is rational on every draw and the
expected work is $[8^kQ_{\rm ex}+3^k]\poly(I)$.
\end{enumerate}
No growth, label-gap, uniqueness, or residual convexity premise is used.
\end{theorem}

For $k=0$ there is no core, and one oracle call with the original bounds
solves the sampled problem exactly. For $k\ge1$ the certificate is the
following. At the incumbent corner $c_j$ with label $z_j$, the $2n_R$
restrictions $z_i\le z_{j,i}-1$ and $z_i\ge z_{j,i}+1$, each together with
the original constraints, have union exactly $Y\setminus\{z_j\}$, however the
residual coordinates are coupled. A restriction that leaves the original
bounds ($z_{j,i}=\ell_i$ in the first case, $z_{j,i}=u_i$ in the second) is
empty and receives the value $+\infty$ without an oracle call, as does a
restriction that the oracle reports infeasible. Let $V_{\rm other}(c_j)$ be
the least of the $2n_R$ values; it is $+\infty$ when all restrictions are
empty, that is, when $z_j$ is the only feasible label; the minimum of an
empty restriction list is also $+\infty$. If
\begin{equation}\label{eq:int:label}
 V_{\rm other}(c_j)-V(c_j)>2G\,w_\infty(\mathcal H_j),
\end{equation}
then $z_j$ is the unique optimal label at every core point of the retained
hull $\mathcal H_j$ (\cref{lem:rec:value}(c)), so every global minimizer lies
in $\mathcal H_j\times\{z_j\}$. The whole slice $[0,1]^k\times\{z_j\}$ is
feasible and contains a global minimizer, so \cref{thm:int:solver} applied to
$F_\gamma(\cdot,z_j)$ on $[0,1]^k$ finishes the problem, whether the core
minimizers are unique, tied, or form a continuum. Feasible values for the
restricted problems would not suffice; \eqref{eq:int:label} needs their exact
optimal values. Distinct integer labels are at unit distance, so point growth
forces \eqref{eq:int:label} at a base-computed level; no label gap is supplied
or charged separately. On the remaining draws the algorithm enumerates the
$N_Z$ labels on the same draw, solves each core slice, and keeps the best by
exact comparisons (\cref{cor:int:mixed-solver}). It returns only the winning
label and its core representation, so the output bound does not depend on
how many labels were examined. The proof is in \cref{app:int:native}.

\begin{corollary}[Implicit core output]\label{cor:int:native-implicit}
Under the hypotheses of \cref{thm:int:native}, an algorithm with expected bit
work $8^kQ_{\rm ex}\poly_d(I)$ returns on ordinary draws the exact label $z^*$
and an implicit patch output for the core slice $F_\gamma(\cdot,z^*)$, and on
the remaining draws the output of \cref{thm:int:native}. A $q$-bit evaluation
of the ordinary output costs $\poly_d(I+q)$, and its expected cost over all
draws is $\poly_d(I+q)$.
\end{corollary}

The variant replaces the exact core solve by the sign and Hessian tests of
\cref{lem:rec:closure} on the core coordinates of the certified slice; it pays
for an additional active-gradient event instead of the term $c_d^k$.

\paragraph{Separable convex recourse over totally unimodular systems.}
The main source of exact integer oracles in this paper is the following
classical result.

\begin{proposition}[Separable convex integer oracles]\label{prop:int:hs}
Let $Y=\{z\in\Z^{n_R}:Az=b,\ \ell\le z\le u\}$, where $A$ is an integer
totally unimodular matrix (promised; automatic for network matrices), $b$ is
integral, and the bounds are integral and finite. Let
$F_0(v,z)=\phi(v)+\sum_if_i(v,z_i)$ with explicit rational polynomials of
fixed degree, and suppose $t\mapsto f_i(v,t)$ is convex on $[\ell_i,u_i]$ for
every $v\in[0,1]^k$ (promised, or certified). Then a box-stable exact integer
oracle exists for every rational linear perturbation, with bit work
polynomial in the query length and logarithmic dependence on the numerical
bounds. This includes integer flows on a directed network with integral
supplies and arc bounds; inequality systems $Cz\le b$ with $C$ totally
unimodular reduce to this form (\cref{cor:int:tu-ineq}).
\end{proposition}

The oracle is the proximity-scaling algorithm of
\citet[Theorem~4.3]{hochbaum1990-convex-separable-optimization-is-not}, which solves separable convex
integer minimization over totally unimodular systems with logarithmically
many linear programs in the numerical bounds and function evaluations at
prescribed points. At a rational core point all evaluations are rational of
polynomial length; tightening coordinate bounds and adding linear terms keep
the class; and convex tangent extensions beyond the native intervals supply
evaluations if the algorithm requests them
(\cref{app:int:native}). The convexity premise concerns each scalar cost on its
whole real interval. For networks, optimality also has a short certificate.

\begin{lemma}[Marginal potentials]\label{lem:int:potentials}
Let $Y$ be the integer flows of a directed network with $s$ nodes, fixed
integral supplies, and arc bounds, where parallel arcs and loops are allowed,
and let each $t\mapsto f_a(t)$ be convex on $[\ell_a,u_a]$. A feasible $z$
minimizes $\sum_af_a(z_a)$ over $Y$ if and only if there are node potentials
$\pi$ for which every arc of the residual network of $z$ has nonnegative
reduced cost, where a forward arc $a=(p,p')$ with $z_a<u_a$ costs
$f_a(z_a+1)-f_a(z_a)$ and a backward arc with $z_a>\ell_a$ costs
$f_a(z_a-1)-f_a(z_a)$, and the reduced cost adds $\pi_p-\pi_{p'}$ along the
residual arc from $p$ to $p'$. If the residual arc costs are rational
numbers of bit length at most $H_m$, shortest-path distances from an added
zero-cost source supply such potentials, of bit length polynomial in $s$ and
$H_m$.
\end{lemma}

This is a certificate, not an algorithm: repeated unit augmentation is not
claimed to run in polynomial time. The characterization holds for arbitrary
real convex costs. The bit-length statement needs rational marginal costs
of polynomial length; explicit rational polynomials of fixed degree have
them at rational core points, while an arbitrary convex real function need
not. With \cref{prop:int:hs}, \cref{thm:int:native} applies to separable
convex integer flows and totally unimodular systems whose costs depend
arbitrarily on the core. For instance,
$F_0(v,z)=\phi(v)+\sum_a\psi_a(z_a)+v^{\mathsf T}Wz$ with convex $\psi_a$ allows
arbitrary core-to-arc coupling $W$; the curvature bound $L$ is that of $\phi$,
and $W$ enters only through its bit length. Convexity of a fixed-degree
univariate $\psi_a$ on its interval can be verified by univariate sign
determination.

\subsection{Core-only noise for integer flows}\label{sec:int:flows}

Under core-only noise the residual costs are not perturbed, and many optimal
labels can persist on every draw. Full point growth then fails, and the
competing-label certificate cannot pass when two labels are optimal at the
optimal core. Throughout this subsection the instance is a core--recourse
instance (\cref{def:rec:instance}) in which $Y$ is the set of integer flows
of a directed network with $s$ nodes and $n_R$ arcs (parallel arcs and loops
allowed), integral supplies, and integral arc bounds $[\ell_a,u_a]$ in
binary; in particular $Y$ is nonempty. Nonemptiness is decided before
sampling by one exact feasibility linear program, whose vertices are integral
by total unimodularity, and an infeasible network is reported as such. The
sampled objective is
\begin{equation}\label{eq:int:flow-model}
 F_\gamma(v,z)=\phi(v)+\sum_af_a(v,z_a)+\gamma^{\mathsf T}v,
 \qquad v\in[0,1]^k,\ z\in Y,
\end{equation}
with explicit rational polynomials of degree at most $d\ge1$ in expanded
form, each $f_a(v,\cdot)$ convex on $[\ell_a,u_a]$ for every $v\in[0,1]^k$
(promised or certified), and a curvature bound $L$ as in \eqref{eq:rec:L}
(verified or certified). Only the $k$ core coefficients are perturbed. By
\cref{prop:model:regret} the sampled optimizer is within $k\sigma$ of
$\min_XF_0$ in original objective value, whatever the size of the network.

A deterministic baseline shows what the noise adds. It uses no flow
structure: exact or certified conditional values on a fixed core grid give
an additive approximation for every core--recourse instance.

\begin{corollary}[A deterministic core grid]\label{cor:int:grid}
Let a core--recourse instance (\cref{def:rec:instance}) with $k\ge1$ and a
fixed rational linear term $\gamma$, for instance $\gamma=0$, be given,
together with a conditional oracle (\cref{sec:rec:value}) that answers
either exactly, in which case put $\eta=0$, or in certified mode with a fixed
accuracy $\eta>0$. For an integer $m\ge1$, query the oracle at all $(m+1)^k$
points of $G_m=\{0,\frac1m,\ldots,1\}^k$, let $(\hat v,\hat z)$ be a returned
completion of least upper value, and put
$\underline f=\min_{v\in G_m}\ell(v)-kL/(8m^2)$. Then $(\hat v,\hat z)\in X$
and
\[
 \underline f\le f^*\le F_\gamma(\hat v,\hat z)
 \le\underline f+\frac{kL}{8m^2}+\eta .
\]
For every $\varepsilon>0$, $m=\lceil\sqrt{kL/(8\varepsilon)}\,\rceil$ gives an additive
gap at most $\varepsilon+\eta$ with $(m+1)^k$ oracle calls. For integer
flows, \cref{prop:int:hs} supplies exact values.
\end{corollary}

The grid rounds only the $k$ core coordinates. The error and the number of
calls do not depend on the number of residual variables, which the oracle
optimizes exactly; rounding them as well would reintroduce a
dimension-dependent error (\cref{ex:rec:star}). The corollary is
deterministic and uses no noise; its proof is in \cref{app:int:flows}. It is
an approximation, not an exact method. With $\sigma=\varepsilon/(2k)$ and an
evaluation of accuracy $\varepsilon/2$, the core-only theorems below also
certify an $\varepsilon$-approximate solution of $F_0$ on every draw
(\cref{prop:model:regret}), but their expected work contains the factor
$[3+(1+k/2)kL/\varepsilon]^k$, whose dependence on $1/\varepsilon$ is worse
than the $(2+\sqrt{kL/(8\varepsilon)})^k$ calls of the grid. The smoothed theorems do not improve additive approximation of
$F_0$; what they add is exact optimization of the sampled objective, with
exact output, on every draw.

A natural exact certificate selects one flow and proves it optimal on the
retained hull by polynomial potentials. Boundary core optima defeat it:
\cref{prop:lim:flow} gives two parallel arcs and a two-dimensional core on
which, with probability $\frac1{16}$ for every grid size, the optimal core
is a vertex with projected growth constant at least one, both flows are
optimal there, and no single flow is optimal on any box at that vertex,
while the core search retains such a box at every level. With $m$ stages,
an algorithm that insists on a single-flow certificate and enumerates
labels when it fails does work at least $2^m/16$ in expectation. The
obstruction concerns that certificate; the instance itself is easy. Two
remedies follow: a premise that all optimal cores are interior, and a
certificate that uses all tied flows at once.

\begin{theorem}[Core-only flow recourse with interior optimal cores]\label{thm:int:flow-interior}
Fix $d$, and assume, in addition to the model \eqref{eq:int:flow-model} with
$k\ge1$, that for every $\gamma\in[-\sigma,\sigma]^k$ every optimal core point
of $F_\gamma$ lies in $(0,1)^k$ (promised, or certified, for instance by
$\partial_iF_0\le-\sigma-\tau'$ on $\{v_i=0\}$ and $\partial_iF_0\ge\sigma+\tau'$
on $\{v_i=1\}$ for all $z$ in the bounding box $\prod_a[\ell_a,u_a]$ and some
rational $\tau'>0$). There is a power of two $M$, computed from the base
instance with $\log_2M\le\poly_d(I)$, such that for core-only noise with
marginal law $U_{\sigma,M}$ the algorithm returns on every draw an optimal
integer flow and an algebraic output of one optimal core and the optimal
value, of length $c_d^k\poly_d(I)$, with $q$-bit refinement in
$c_d^k\poly_d(I+q)$, and with expected bit work
$[8^kQ_{\rm ex}+c_d^k]\poly_d(I)$. For quadratic objectives the output is
rational.
\end{theorem}

The interiority premise is used only in the running-time analysis;
correctness holds on every draw without it, and its sufficient certificate
enters no numerical factor. The premise is not vacuous because $Y$ is
nonempty. At the incumbent corner the algorithm computes an optimal flow, a
tight shortest-path tree of its residual network, and the corresponding
polynomial potentials, and it tests exactly whether every residual reduced
cost is nonnegative on the whole retained hull; identically zero reduced
costs pass, so persistent ties are allowed. If the test passes, the flow is
optimal on the hull, and the exact solver finishes its whole core slice. The
analysis uses a finite family of reduced-cost polynomials over all labels
and trees, fixed before sampling: an optimal core in the interior is a
stationary point of every attaining fixed-flow slice, so the noise equals a
negative gradient there, and the event that the optimal core is near a zero
of a chart polynomial has small probability by an algebraic tube estimate.

\paragraph{All tied flows at once.}
At a core point $c$ on a face $\mathcal A$ of the core cube, let $z^0$ be an
optimal flow and $\pi$ potentials certifying it (\cref{lem:int:potentials}).
The adjusted costs $g_a(c,t)=f_a(c,t)+(\pi_p-\pi_{p'})t$ of the arcs $a=(p,p')$
are convex, and their integer minimizers form intervals $I_a\ni z^0_a$, found
by binary search on monotone first differences. The set
$Y_0=\{z\in Y:z_a\in I_a\text{ for all }a\}$ is exactly the set of optimal flows
at $c$: the potential terms sum to a constant over $Y$, and the objective gap is
a sum of nonnegative scalar gaps. Minimizing an inward core derivative over
$Y_0$ is again a convex flow problem, because on an interval with at least
three integers the adjusted cost is constant, so the original arc cost is
affine there and its inward derivative is convex; two-point intervals use
linear interpolation, and singleton intervals are substituted. Finally, every
feasible flow is close to $Y_0$.

\begin{lemma}[Proximity to the optimal set]\label{lem:int:proximity}
For every $z\in Y$ there is $\bar z\in Y_0$ with
$\norm{z-\bar z}_1\le n_R\sum_a\dist(z_a,I_a)$.
\end{lemma}

The deterministic certificate uses these facts on a rational core box
$\mathcal D$ whose intervals in the normal coordinates of $\mathcal A$ contain
the corresponding bounds, with projection $\mathcal D_{\mathcal A}$ onto the
face and $c\in\mathcal D_{\mathcal A}$. Let $\varsigma_i=+1$ for a normal
coordinate fixed at $0$ and $\varsigma_i=-1$ for one fixed at $1$, let $T_1$ and
$T_\infty$ be the sum and the maximum of the normal widths of $\mathcal D$, let
$R_{\mathcal A}$ bound $\norm{w-c}$ on $\mathcal D_{\mathcal A}$, and let
$H\ge1$ and $K\ge0$ be rational bounds on $\norm{\nabla^2_vF_0}$ and on
$\abs{\partial_{v_i}\partial_{z_a}f_a}$ on the real bounding domain.

\begin{theorem}[An optimal-face certificate]\label{thm:int:face}
Suppose that, with the potentials taken as polynomials of the core obtained
from a tight shortest-path tree at $c$:
\begin{enumerate}[label=(\alph*)]
\item for every arc, $g_a(w,t)=g_a(w,z^0_a)$ for all $w\in\mathcal D_{\mathcal A}$
and all integers $t\in I_a$;
\item every first adjusted marginal outside $I_a$ that stays in the native
interval is at least $\mu_{\rm out}=n_RKT_1$ on $\mathcal D_{\mathcal A}$; and
\item if $\mathcal A$ has normal coordinates, $\beta=\min_i\beta_i$ satisfies
$\beta-HR_{\mathcal A}-\frac H2T_\infty>0$, where
$\beta_i=\min_{z\in Y_0}\varsigma_i\partial_{v_i}F_\gamma(c,z)$.
\end{enumerate}
Then every flow of $Y_0$ is optimal at every point of $\mathcal D_{\mathcal A}$,
and every point of $\mathcal D\times Y$ off the face is strictly worse than some
feasible point on it. If $\mathcal D$ contains the core of every global
minimizer, the slice $[0,1]^k\times\{z^0\}$ contains a global minimizer.
Conditions (a)--(c) are decided exactly, (a) by at most $d+1$ polynomial
identities per arc, (b) by exact polynomial minimization on a rational box
(\cref{thm:int:solver}), and (c) by one convex flow problem per normal
coordinate.
\end{theorem}

The certificate is sound on every draw and needs neither uniqueness of the
optimal flow nor a noise assumption. Both (b) and (c) are needed: (c) uses the
minimum over \emph{all} optimal flows, since a single tied flow can have the
wrong inward sign, and without (b) a flow that loses at $c$ can become optimal
off the face. In the instance of \cref{prop:lim:flow} the certificate with the origin
face closes every box of maximum width less than $2\min_i\gamma_i$.

\begin{theorem}[Core-only flow recourse at arbitrary core faces]\label{thm:int:flow-boundary}
Fix $d$. For the model \eqref{eq:int:flow-model} with $k\ge1$ there are an
effective function $f_d$ and a power of two $M$, computed from the base
instance with
\[
 \log_2M\le f_d(k)\poly_d(I),
\]
such that for core-only noise with marginal law $U_{\sigma,M}$ the algorithm
returns on every draw an optimal integer flow and an algebraic output of one
optimal core and the optimal value, of length $f_d(k)\poly_d(I)$, with
$q$-bit refinement in $f_d(k)\poly_d(I+q)$, and with expected bit work
$f_d(k)Q_{\rm ex}\poly_d(I)$. No interiority, growth, or uniqueness premise is
used, and arbitrarily many optimal flows may persist.
\end{theorem}

At each level the algorithm tries the certificate of \cref{thm:int:face} on
every face of the core cube compatible with the retained hull, at the midpoint
of the projected hull, and solves the whole core slice of $z^0$ when one passes.
The analysis needs three probability bounds: projected growth, distance of the
optimal core from the zeros of a base-only family of chart polynomials (tube
estimate), and a margin $\tau$ for the minimal inward derivative over all
optimal flows. The last is obtained without conditioning on a random label: on
each fixed face, the optimal core and its whole optimal flow set do not depend
on the normal noise coefficients. A tube controls the distance from a zero set,
not the value of a polynomial away from it. Converting distance into a value
margin uses fixed-block quantifier elimination with its coefficient-height
bound, and costs a precision of $f_d(k)\poly_d(I)$ bits. This is why the
sampling precision in \cref{thm:int:flow-boundary} depends on $k$.

\begin{corollary}[Bilinear coupling]\label{cor:int:bilinear}
If $F_\gamma(v,z)=\phi(v)+\sum_a\psi_a(z_a)+v^{\mathsf T}Wz+\gamma^{\mathsf T}v$
with a rational matrix $W$, fixed degree $d\ge2$, convex $\psi_a$ on
$[\ell_a,u_a]$, and $\partial_{ii}\phi\le L$ on $[0,1]^k$, then the conclusions
of \cref{thm:int:flow-boundary} hold with $\log_2M\le\poly_d(I)$, output length
$c_d^k\poly_d(I)$, refinement $c_d^k\poly_d(I+q)$, and expected bit work
$[8^kQ_{\rm ex}+c_d^k]\poly_d(I)$. For quadratic $\phi$ the output is rational.
\end{corollary}

All chart polynomials are then affine in the core, and their value margins
have polynomial length. Let $p$ be such a chart, nonzero, with coefficients
of bit length at most $H_0$, on a face identified with $[0,1]^q$. If its zero
set $Z$ in the cube is nonempty, then $p$ is not constant and
$\abs{p(x)}\ge b_{\min}\dist(x,Z)$, where $b_{\min}\ge2^{-H_0}$ is the least
absolute value of a nonzero linear coefficient. If $p$ is a nonzero constant
or has no zero in the cube, then $\abs p$ is at least its least absolute
vertex value, a positive rational at least $2^{-(q+1)H_0}$; this case includes
vertex faces. Identically zero charts are discarded from the exceptional
family, and they pass the non-strict certificate tests. The curvature in the
numerical factor is that of $\phi$ alone. This would fail for nonlinear
coupling: a term $v_i^2z_a^2$ makes the core curvature grow quadratically with
the capacity of arc $a$.

\subsection{Totally unimodular systems}\label{sec:int:tu}

\begin{theorem}[Totally unimodular recourse]\label{thm:int:tu}
Fix $d$ and $k\ge1$. Let $Y=\{z\in\Z^{n_R}:Az=b,\ \ell\le z\le u\}$ be
nonempty, where $A$ is an integer totally unimodular matrix (promised), $b$
is integral, and the bounds are finite integers encoded in binary. Let
\[
 F_\gamma(v,z)=\phi(v)+\sum_if_i(v,z_i)+\gamma^{\mathsf T}v
 \qquad(v\in[0,1]^k,\ z\in Y),
\]
where $\phi$ and the $f_i$ are explicit rational polynomials of degree at most
$d$ in expanded form, each $f_i(v,\cdot)$ is convex on $[\ell_i,u_i]$ for
every $v\in[0,1]^k$ (promised or certified), $L>0$ is a rational bound
\eqref{eq:rec:L} on $[0,1]^k\times\prod_i[\ell_i,u_i]$ (verified or
certified), and only the $k$ core coefficients are perturbed, with law
$U_{\sigma,M}$. Then the conclusions of \cref{thm:int:flow-boundary} hold. If
moreover $F_\gamma=\phi(v)+\sum_i\psi_i(z_i)+v^{\mathsf T}Wz+\gamma^{\mathsf T}v$
with the premises of \cref{cor:int:bilinear}, its conclusions hold.
\end{theorem}

These are the premises of the flow model, with the network replaced by a
totally unimodular system; certificates are part of the input length, and
nonemptiness is decided by one linear program. The premises are used: the
proof evaluates and differentiates explicit rational polynomials at rational
core points, needs rational unit marginals of polynomial length for its
charts, and uses $L$ for the count. Total unimodularity with arbitrary convex
costs, or with succinctly encoded costs, does not provide these interfaces.
Two further ingredients replace network structure. Shortest-path potentials
become a dual vector satisfying at most $2n_R$ adjacent-slope inequalities,
which exist because the piecewise-linear interpolation of the separable costs
has an integral optimum over the real totally unimodular polytope; a dual
vertex in a pre-draw box has inverse-basis entries in $\{0,\pm1\}$, which
bounds the chart coefficients. Directed cycles become conformal circuits of
$\ker A$ with entries in $\{0,\pm1\}$, which gives \cref{lem:int:proximity}
with the same factor $n_R$. Both ingredients use total unimodularity; neither
holds for general integer matrices.

\begin{corollary}[Bounded totally unimodular inequalities]\label{cor:int:tu-ineq}
Under the premises of \cref{thm:int:tu} on the core and the costs, its
conclusions hold for a nonempty
$Y=\{z\in\Z^{n_R}:Cz\le b,\ \ell\le z\le u\}$ with $C$ totally unimodular, $b$
integral, and finite integral bounds.
\end{corollary}

Add the slack $s_i=b_i-C_iz$ with the explicit bounds
$0\le s_i\le b_i-\min_{\ell\le z\le u}C_iz$; the matrix $[C\ I]$ is totally
unimodular, projection is a bijection onto $Y$, and slacks have zero cost, so
the curvature bound and all objective values are unchanged. A negative slack
bound proves infeasibility.

\subsection{Strong noise and random components}\label{sec:int:strong-field}

The previous results allow every noise scale $\sigma>0$, with numerical
ratios such as $L/\sigma$ in their bounds, and they need a small core. The
following result needs neither a core nor a decomposition, but it requires a
lower bound on the noise: noise that dominates the variation of every
coordinate derivative. Let $X=\prod_iX_i$ be a mixed box
(\cref{sec:model:input}) with $n$ coordinates after preprocessing, and let
$F_0$ be an explicit rational polynomial of degree at most $d$. If $n=0$ the
objective is a constant, which is returned directly; assume $n\ge1$. The
\emph{primal graph} joins two coordinates when they occur together in a
nonzero monomial, so each monomial support is a clique; let
$\Delta_+=\max\{1,\text{maximum degree}\}$. Before sampling, compute rational
enclosures $\partial_iF_0\in[\underline\partial_i,\overline\partial_i]$ on the
continuous hull, by multiplying exact interval ranges of univariate powers
over each derivative monomial. Put $R_i=\overline\partial_i-\underline\partial_i$,
\[
 a_i=\begin{cases}c_d,&i\text{ continuous},\\u_i-\ell_i+1,&i\text{ native integer},
 \end{cases}\qquad
 q_i=U_{\sigma,M}\bigl([-\overline\partial_i,-\underline\partial_i]\bigr),\qquad
 \beta=\max_ia_iq_i,
\]
where $c_d$ is the constant of \cref{thm:int:solver}; for quadratic
objectives one may take $a_i=3$ for continuous coordinates. Each $q_i$ is an
exactly computable rational number, which depends on $M$.

\begin{theorem}[Strong fields and random components]\label{thm:int:strong-field}
Fix $d$, and let $M\ge2$ be any power of two. Under the ambient perturbation
model with marginal law $U_{\sigma,M}$, a deterministic algorithm returns on
every draw an exact global minimizer of $F_\gamma$ on $X$ in the component
format of \cref{def:model:outputs}\ref{it:model:component}: rational values
for the pinned coordinates, and for each remaining component its integer labels
and one common-root algebraic representation of its continuous coordinates and
value. If $4\Delta_+\beta<1$, its expected bit work is at most
\begin{equation}\label{eq:int:sf-work}
 \poly_d(I+\log_2M)\Bigl[1+\frac{\sum_ia_iq_i}{1-4\Delta_+\beta}\Bigr].
\end{equation}
A $q$-bit evaluation of the output, that is, a rational point within
Euclidean distance $2^{-q}$ of the represented optimizer and a rational
interval of width at most $2^{-q}$ containing the optimal value, together with
a feasible rational point whose objective gap is certified to be at most
$2^{-q}$, has expected cost bounded by \eqref{eq:int:sf-work} with
$\poly_d(I+\log_2M+q)$ in place of $\poly_d(I+\log_2M)$. For quadratic
objectives the output is an exact rational optimizer and value, with $a_i=3$
for continuous coordinates. A sufficient condition for a bounded denominator
is
\begin{equation}\label{eq:int:sf-regime}
 \sigma\ge8\Delta_+\max_i(a_iR_i),\qquad M\ge16\Delta_+\max_ia_i,
\end{equation}
which gives $a_iq_i\le1/(8\Delta_+)$ and $4\Delta_+\beta\le\frac12$.
\end{theorem}

Outside the closed interval in the definition of $q_i$, the sampled
derivative $\partial_iF_\gamma$ has a strict sign throughout the continuous
hull, so every global minimizer has $x_i$ at the corresponding bound; this
holds for integer coordinates as well, by monotonicity between labels. Call
a coordinate \emph{bad} when its coefficient lies inside that closed
interval, so this sign test leaves it unpinned. These bad events depend on
one coefficient each, so they are independent
Bernoulli events with the exact probabilities $q_i$; recomputing the
enclosures after pinning would destroy this independence and is not needed.
After pinning, every surviving monomial lies in one connected component of the
graph induced by the bad coordinates, so the objective splits into independent
component problems, each solved exactly by \cref{cor:int:mixed-solver} or, for
quadratics, by face enumeration. At most $(4\Delta_+)^{t-1}$ connected sets of
size $t$ contain a given vertex, and a subcritical branching bound gives
\eqref{eq:int:sf-work}. The $q$-bit evaluation refines each of the at most $n$
components with $q+\lceil\log_2(n+1)\rceil$ bits, for its point, its value,
and its certified objective gap separately. Combining their point errors
in Euclidean norm and adding their value widths and gaps gives the stated
joint contract; an objective gap alone would not bound the distance to the
optimizer on a flat or singular component. The proof is in
\cref{app:int:strong-field}.

The algorithm is exact on every draw, also when $4\Delta_+\beta\ge1$, but the
theorem then gives no useful bound. The probabilities $q_i$ are not monotone
in $M$: grids of sizes $M$ and $2M$ are not nested, so an interval without
atoms of one grid can contain atoms of the other. A value of $\beta$
computed for one grid size therefore does not transfer to another; $q_i$ and
$\beta$ must be computed for the $M$ actually used. The sufficient condition
\eqref{eq:int:sf-regime}, in contrast, holds for every larger power of two.
This sufficient regime scales with the full derivative
variation, with the graph degree, and, for integer coordinates, with the
numerical number of labels; for continuous coordinates of nonquadratic
objectives it scales with the large effective constant $c_d$. The total value
is reported as a rational constant plus a sum of component values. Expanding
that sum into one defining polynomial can have degree exponential in the
number of components, and exact sign or equality tests for sums of unrelated
algebraic numbers are not claimed; rational enclosures of the total are
obtained by refining each component. The regret bound of
\cref{prop:model:regret} is $\sigma\sum_iw_i$, which is large in this regime.

\subsection{Separable integer problems with a low-rank concave part}\label{sec:int:lowrank}

The last result uses aligned noise (\cref{def:model:perturbation}) and an exact
argument of a different kind: the original objective values of the sampled
problem lie on a lattice whose spacing is controlled by the same grid that
controls the expected count. Let $X=\prod_{j=1}^nX_j$ with
$X_j=\{L_j,\ldots,U_j\}$ and integer bounds $L_j\le U_j$ in binary, and
\begin{equation}\label{eq:int:lowrank-model}
 F(x)=\sum_{j=1}^ng_j(x_j)-\frac\alpha2\norm{Tx}^2,
\end{equation}
with a rational $\alpha>0$ and an arbitrary supplied $T\in\Q^{k\times n}$,
whose rows may be zero, linearly dependent, or constant on $X$. Fix a degree
bound $d_*$. Each $g_j$ is a continuous convex function on $[L_j,U_j]$, given
by an explicit list of polynomial pieces of degree at most $d_*$, in expanded
monomial form with rational coefficients, on consecutive intervals with
rational breakpoints that cover $[L_j,U_j]$; a single polynomial is a list with
one piece. Fixed-degree shifted or factored pieces can first be expanded in
polynomial time; their expanded coefficients define the denominator budget.
Convexity is verified: continuity and the order of the one-sided
derivatives at each breakpoint are rational evaluations, and nonnegativity of
the second derivative of each piece on its interval is a univariate sign
determination. Smoothness at the breakpoints is not assumed. For $k\ge1$
choose a rational $\sigma_i>0$ for each row, let $\ell_i=\min_X(Tx)_i$,
$u_i=\max_X(Tx)_i$, $w_i=u_i-\ell_i+2\sigma_i/\alpha$, and $s=\max_iw_i$. Let
$Q_{\rm den}$ be the product of the denominators of all base rational data,
with the pieces in expanded form, put $D_0=2Q_{\rm den}^3$, and let $M$ be the
least power of two with $M\ge\max\{2,k\alpha s^2D_0\}$. The sampled objective
is $F_\xi(x)=F(x)+\xi^{\mathsf T}Tx$, where $\xi_1,\ldots,\xi_k$ are
independent and $\xi_i$ has law $U_{\sigma_i,M}$; all rows use the same $M$.

\begin{theorem}[Exact low-rank separable integer optimization]\label{thm:int:lowrank}
If $k=0$, a deterministic algorithm minimizes $F$ on $X$ exactly in
polynomial time. If $k\ge1$, then for every draw a deterministic algorithm
returns an exact integer minimizer of $F_\xi$ on $X$ and the exact rational
optimal value, both of polynomial length. Its expected bit work is at most
\[
 8^kH_{\rm rat}P(I),\qquad
 H_{\rm rat}=\prod_{i=1}^k\Bigl[3+\frac{(1+2k)\alpha w_i}{2\sigma_i}\Bigr],
\]
where $P$ is a polynomial with an absolute exponent and coefficients
depending only on $d_*$. No growth, uniqueness, or separation premise is
used, and there is no fallback. The number of integer coordinates, the
numbers of pieces, and the numerical lengths of the integer intervals do not
enter $H_{\rm rat}$ except through the row ranges $u_i-\ell_i$.
\end{theorem}

Completing the square gives the auxiliary value
$W(a)=\frac\alpha2\norm a^2+\min_{x\in X}[G(x)-\alpha a^{\mathsf T}Tx]$ on the
fixed box $\prod_i[\ell_i-\sigma_i/\alpha,u_i+\sigma_i/\alpha]$, with
$G=\sum_jg_j$. It has upper coordinate curvature $\alpha$, because
$W-\frac\alpha2\norm a^2$ is a minimum of affine functions. Its inner problem
separates into scalar problems over consecutive integers. Real convexity of
$g_j$ makes the forward differences of $g_j(t)-\lambda t$ nondecreasing, also
between different pieces, so binary search on them finds an exact integer
minimizer with $O(1+\log(U_j-L_j+1))$ evaluations, without listing labels.
Moreover $\min_aW(a)+\xi^{\mathsf T}a$ equals $\min_XF_\xi$ up to the known
constant $\norm\xi^2/(2\alpha)$. The corrected-corner search of
\cref{thm:count:cells} on the $k$ auxiliary coordinates gives an integer
witness whose original gap is below $E_J<1/(D_0(M-1))$ at the predetermined
level $J=\log_2M$, while every original value $F_\xi(x)$ lies in
$\frac1{D_0(M-1)}\Z$. At an integer argument every expanded monomial is an
integer, so the unary degree and the pieces add no denominator, and one
common grid denominator suffices because the noise enters the original
objective linearly. The auxiliary values may have denominators involving
$(M-1)^2$, but no lattice property of them is used. The proof is in
\cref{app:int:lowrank}.

The argument uses discrete convexity of the unaries, their exact evaluation,
and a common denominator of their values, and nothing else about them. The
fixed degree and the explicit piece lists guarantee exact evaluation in
polynomial time with values of polynomial length. A succinct encoding would
not: evaluating a monomial $t^{2^b}$ with a binary exponent at $t=2$
produces an integer of length $2^b$. Continuous coordinates and coupled
integer constraints are outside the theorem: continuous minima can be
irrational, and coupling destroys the separable scalar recourse.

The numerical ratios $\alpha w_i/\sigma_i$ are explicit; binary-encoded bounds
alone do not make them small. For binary variables, and more generally for
explicitly listed labels, deterministic fixed-rank algorithms already exist. On
$\{0,1\}^n$ each $g_j$ agrees with the affine function $g_j(0)+c_jx_j$,
$c_j=g_j(1)-g_j(0)$, so $F$ is a concave
function of $(Tx,c^{\mathsf T}x)$; it attains its minimum at a vertex of the
zonotope that is the image of the cube, and this zonotope in dimension $k+1$
has $O(n^k)$ vertices. With explicitly listed labels, the planar polytopes
$\conv\{(T_{\cdot j}t,g_j(t)):t\in X_j\}$ give a Minkowski-sum enumeration
that, for fixed $k$, is polynomial in the total number of labels. The theorem
does not improve these cases; it allows binary-encoded intervals without listing
their labels.

\subsection{Scope}\label{sec:int:scope}

The residual feasible set must not depend on the core, and the exact integer
oracle must handle all tightened coordinate bounds. The flow and totally
unimodular results need nonempty residual sets and separable costs that are
convex in each residual coordinate for every core point; they do not cover
nonseparable integer convexity, arbitrary integer matrices, or core-dependent
supplies and right-hand sides. Under core-only noise, single-flow certificates
fail at boundary cores (\cref{prop:lim:flow}); the face certificate needs a
parameter-dependent sampling precision for nonlinear coupling, which bilinear
coupling avoids. The strong-field theorem is a strong-noise statement, and the
low-rank theorem uses aligned rather than ambient noise. All results optimize
the sampled objective exactly; none optimizes $F_0$.

The native, interior-flow, and bilinear results, including the bilinear
totally unimodular case, and the lattice theorem have polynomial sampling
precision. By \cref{cor:count:universal-law}, with a fixed oracle
implementation where an oracle is used, their grid size can therefore be
replaced by one depending only on the input length; in the lattice theorem
the level cap is reset to $\log_2M$. For the nonlinear results at arbitrary
core faces (\cref{thm:int:flow-boundary} and the corresponding case of
\cref{thm:int:tu}), the level cap and the grid exponent are bounded by
$F_d(k)P_d(I)$ for an effective nondecreasing $F_d$. A common resolution then
exists for the instances with a supplied bound $K\ge k$, and $F_d(K)$ enters
their sampling, output, and work bounds; taking $K=I$ gives a law that
depends only on $I$, but its sampling cost need not stay within
$f_d(k)\poly_d(I)$ for the actual $k$. The strong-noise theorem holds for
every $M$, with $q_i$ and $\beta$ computed for the $M$ used.

The integer oracles, fixed-block elimination, tube estimates, and
connected-set counting are established tools; what this section proves is
their composition with one base-selected finite law, the label and face
certificates, and the bit and output bounds, including exact algebraic output
in which all coordinates and the value within one completion or component
share one root.
